% SPDX-License-Identifier: CC-BY-4.0
\documentclass[11pt,letterpaper]{article}

\usepackage[T1]{fontenc}
\usepackage{amsmath,amssymb}
\usepackage{tikz}
\usepackage{pgfplots}
\pgfplotsset{compat=1.18}
\usepackage{algorithm}
\usepackage{algpseudocode}
\usepackage{booktabs}
\usepackage{enumitem}
\usepackage[hidelinks]{hyperref}

\newcommand{\N}{\mathbb{N}}
\newcommand{\Z}{\mathbb{Z}}
\newcommand{\Q}{\mathbb{Q}}
\newcommand{\R}{\mathbb{R}}
\newcommand{\C}{\mathbb{C}}
\newcommand{\Map}[3]{#1 \colon #2 \to #3}
\newcommand{\Mod}[1]{#1\text{-}\mathrm{Mod}}
\newcommand{\GroupHom}{\mathrm{Hom}_{\mathbf{Grp}}}
\newcommand{\RingHom}{\mathrm{Hom}_{\mathbf{Ring}}}

\newtheorem{theorem}{Theorem}
\newtheorem{lemma}[theorem]{Lemma}
\newtheorem{corollary}[theorem]{Corollary}
\newtheorem{definition}[theorem]{Definition}
\newtheorem{remark}[theorem]{Remark}
\newenvironment{proof}{\noindent\textit{Proof.}}{\hfill\(\square\)}

\title{Generalized Collatz Library}
\author{}
\date{\today}

\begin{document}
\maketitle
\section*{Common-power time compression}

\section{Common-power time compression}

% The signed divisor and the raw branch order remain part of the definition.
\begin{definition}
For integers \(q,c,d\) with \(|d|\ge2\), let
\[
F_{q,c,d}(n)=
\begin{cases}
n/d,&d\mid n,\\
qn+c,&d\nmid n.
\end{cases}
\]
Suppose that, for some \(a\in\N\) and \(Q,C\in\Z\),
\[
q=d^aQ,
\qquad
c=d^aC.
\]
Write \(F=F_{q,c,d}\) for the original map and \(G=F_{Q,C,d}\) for the
compressed map. Define the state-dependent block length
\[
\tau(n)=
\begin{cases}
1,&d\mid n,\\
a+1,&d\nmid n.
\end{cases}
\]
Both maps act on all of \(\Z\), and the divisor retains its sign.
\end{definition}

% The forced divisions are proved directly and do not use an accelerated map.
\begin{lemma}[QCD.TimeCompression.001.1a]
For every \(b\in\N\) and \(x\in\Z\),
\[
F^b(d^b x)=x.
\]
\end{lemma}
\begin{proof}
Induct on \(b\). The case \(b=0\) is the identity. If the assertion holds for
\(b\), then
\[
d^{b+1}x=d(d^b x)
\]
is divisible by \(d\), including when \(d<0\). The first raw instruction is
therefore division and sends this state to \(d^b x\). The next \(b\)
instructions send it to \(x\) by the induction hypothesis.
\end{proof}

% Each branch is compared at its actual raw time rather than at one artificial step.
\begin{theorem}[QCD.TimeCompression.001.1]
For every \(n\in\Z\),
\[
F^{\tau(n)}(n)=G(n).
\]
The clock is positive: \(\tau(n)\ge1\).
\end{theorem}
\begin{proof}
If \(d\mid n\), both maps take the division branch, and \(\tau(n)=1\), so both
values are \(n/d\).

Suppose \(d\nmid n\). The original affine instruction gives
\[
F(n)=qn+c=d^a(Qn+C).
\]
The factor \(d^a\) makes the following \(a\) original instructions forced
divisions. Lemma~QCD.TimeCompression.001.1a sends their starting state to
\(Qn+C\). Since the compressed guard is the same and \(d\nmid n\), this is
exactly \(G(n)\). Here \(\tau(n)=a+1\). Both possible values of \(\tau\) are at
least one.
\end{proof}

\section{The complete sampled trajectory}

% The event count uses precisely the states at which a compressed instruction starts.
\begin{definition}
For \(k\in\N\), let
\[
R(k,n)=\#\{j:0\le j<k,\ d\nmid G^j(n)\}
\]
be the number of affine instructions among the first \(k\) compressed steps.
Let
\[
T(k,n)=\sum_{j=0}^{k-1}\tau(G^j(n))
\]
be the accumulated original raw time, with both quantities zero when \(k=0\).
\end{definition}

% An affine block has one sampled instruction and \(a\) extra divisions.
\begin{lemma}[QCD.TimeCompression.001.2a]
For every \(k\in\N\) and \(n\in\Z\),
\[
T(k,n)=k+aR(k,n).
\]
In particular, \(T(k,n)\ge k\).
\end{lemma}
\begin{proof}
Each of the \(k-R(k,n)\) division instructions contributes one to the sum. Each
of the \(R(k,n)\) affine instructions contributes \(a+1\). Therefore
\[
T(k,n)=(k-R(k,n))+(a+1)R(k,n)=k+aR(k,n).
\]
The final inequality follows because \(aR(k,n)\ge0\).
\end{proof}

% Induction composes blocks in chronological order and accounts for every sampled state.
\begin{theorem}[QCD.TimeCompression.001.2]
For every \(k\in\N\) and \(n\in\Z\),
\[
G^k(n)=F^{T(k,n)}(n)=F^{\,k+aR(k,n)}(n).
\]
\end{theorem}
\begin{proof}
Induct on \(k\). At \(k=0\), both iterates equal \(n\), and \(T(0,n)=0\).
Assume the first equality at time \(k\). Splitting the original clock after the
first \(T(k,n)\) steps and applying the one-block theorem at \(G^k(n)\) gives
\[
F^{\,T(k,n)+\tau(G^k(n))}(n)
 =F^{\tau(G^k(n))}(G^k(n))
 =G^{k+1}(n).
\]
The exponent on the left is \(T(k+1,n)\) by definition. This proves the first
equality at \(k+1\); the clock formula proves the second.
\end{proof}

\section{Recovery and limitations}

% The return statement concerns the exact expanded time and makes no minimality inference.
\begin{theorem}[QCD.TimeCompression.001.3]
For every \(k\in\N\) and \(n\in\Z\),
\[
F^{\,k+aR(k,n)}(n)=n
\quad\Longleftrightarrow\quad
G^k(n)=n.
\]
Thus a closed compressed itinerary of length \(k\), containing \(r\) affine
events, expands to a closed original raw itinerary of length \(k+ar\).
\end{theorem}
\begin{proof}
At the displayed original time, the sampled-trajectory theorem identifies the
left-hand iterate with \(G^k(n)\). Substitution gives the equivalence. If the
itinerary has \(r=R(k,n)\) affine events, the same identity gives the stated
length.
\end{proof}

% Arbitrary targets need no invariance because the compressed endpoint is an actual raw state.
\begin{theorem}[QCD.TimeCompression.001.4]
Let \(S\subseteq\Z\) be arbitrary. If
\[
\forall n\in\Z\ \exists k\in\N:\ G^k(n)\in S,
\]
then
\[
\forall n\in\Z\ \exists t\in\N:\ F^t(n)\in S.
\]
For a witness \(k\) in the first assertion, one may take
\(t=k+aR(k,n)\) in the second.
\end{theorem}
\begin{proof}
Fix \(n\), and choose \(k\) such that \(G^k(n)\in S\). The sampled-trajectory
theorem gives
\[
F^{\,k+aR(k,n)}(n)=G^k(n)\in S.
\]
This uses the same initial state and the same target set, so no target invariance
or reconstruction of a different seed is required.
\end{proof}

Define the intrinsic coefficient height by
\[
h(u,v)=\max\{|u|,|v|\}.
\]
The factorization gives the exact natural-number identity
\[
h(q,c)=|d|^a h(Q,C).
\]
If \(a>0\) and \((Q,C)\ne(0,0)\), then \(|d|^a>1\) and \(h(Q,C)>0\), so
\[
h(Q,C)<h(q,c).
\]
Thus the replacement is a strict reduction in coefficient height precisely in
this nondegenerate case. The divisor, signed state domain, guards, all initial
states, and every sampled state are retained. If \((Q,C)=(0,0)\), the original
and compressed coefficient pairs are both zero; the theorem remains exact
transport but supplies no strict decrease. Recovery is the identity of
Theorem~QCD.TimeCompression.001.2. The new obligation is the corresponding target
statement for \(G\); the theorem does not assume that obligation is true.

The identity state correspondence need not intertwine the maps one step at a
time. For example,
\[
F_{4,4,2}(1)=8\ne2=F_{1,1,2}(1),
\]
but the full expanded block is
\[
1\longmapsto8\longmapsto4\longmapsto2,
\qquad
F_{4,4,2}^{\,3}(1)=F_{1,1,2}(1).
\]
The return theorem states the exact expanded return time. It does not assert that
this return is primitive, nor does forward target transport supply a converse for
an arbitrary target that may be met only inside a skipped block.

The Go sweep compares every bounded block on the grid
\[
\begin{gathered}
q,c\in\{-2,-1,0,1,2\},\qquad d\in\{-4,-3,-2,2,3,4\},\\
0\le a\le2,\qquad -15\le n\le15,\qquad 0\le k\le6,
\end{gathered}
\]
checking endpoints, raw step counts, affine event counts and the clock identity.
It produces 97{,}650 certificates.
Arithmetic beyond fixed-width integers and the rejected-parameter boundaries
have separate tests. The finite sweep illustrates the identities above; it does
not replace their proofs.

\end{document}
